We study whether a robot that keeps an explicit, typed memory of the hidden physics of the things it
handles (mass, friction, crush limit, clearance) makes better manipulation decisions than one that does not.
We build (i) a typed, versioned, causal property-graph memory with provenance and an actor/evaluator
boundary, (ii) a 2.28\,M-parameter world model whose interaction history reaches the outcome head only through
readable, overwritable typed slots, (iii) uncertainty-aware planners that deviate from a scripted skill only
when an ensemble lower bound supports it, and (iv) an adaptation of the pretrained looped transformer
\inkling{} into a slow \sysII{} (the unmodified language model) and a fast \sysI{} (gated LoRA and memory
cross-attention active only after a mode token).
Everything is evaluated in simulation, on procedural MuJoCo industrial stations with UR5e, Panda and a held-out
Kinova Gen3 arm, a gantry, and a Unitree G1, using paired tasks, $n\ge 30$ per cell, 95\% bootstrap confidence
intervals, thresholds tuned on validation only, and pre-registration of the later rounds.
The results are mixed and we report them as such.
A typed friction slot trained in a first sprint was \emph{not} causally used.
In a diagnostic case a redesigned model sharpens its fragility estimate after one probe and carries it across a lot; it beats a
generic-memory model on all four robots, but it does not beat a simple probe-then-act script on any arm.
Its one pre-registered win over that script is on the gantry under a shared call budget:
$+0.09$ \ci{+0.03}{+0.16} good parts per lot, attributable to lot memory ($+0.17$ \ci{+0.11}{+0.23} against the
memory-off ablation).
Absolute success on the arms is far from industrial levels and is bounded by the scripted skill layer.
Comparisons against end-to-end imitation (ACT) and the closed-loop \sysI/\sysII{} evaluation are pre-registered
and pending.
